\section{Hearing the texts in the order of Mass}

\subsection{The voice that opens the Mass}

A single person speaks in the Entrance antiphon. Psalm 25 (24) begins with a soul lifted to God and immediately mentions enemies. Someone else will be there to laugh if the confidence proves misplaced. The prospect of shame gives the prayer a public cost. Even before the psalmist asks what road to follow, he has put himself in the position of someone whose reliance on God can be seen. By verse 3 he has widened his claim to include everyone who waits for the Lord. A whole assembly can take up that singular voice. Someone who can barely bring herself to pray and someone who has been waiting a lifetime can say the same words.

Augustine understands Christ to be speaking in the person of the Church. He imagines what has brought her to the point of speaking at all. Earthly desires had drawn her down. There was a time when she trusted her own power and had reason to be ashamed of the outcome. Those who encouraged her then can turn around and make sport of her now. The appeal to God includes a renunciation of the reliance that put her in their hands. Bellarmine begins from the psalmist's reason for trusting. The people who have waited for God before him give him ground to expect that God's promise can bear the weight he places on it. The endangered individual can appeal to what is known of God's dealings with others. Both authors take the move from one person to many seriously.\footnote{Augustine, \textit{Exposition on Psalm} 24, §§1--3 (English heading Psalm 25). Robert Bellarmine, \textit{A Commentary on the Book of Psalms}, Psalm 24, explanations 1--3, trans. John O'Sullivan.}

At the Collect the priest prays for the whole people. The meeting he hopes for has a destination, a place at Christ's right hand. Arriving there involves deeds that can be judged. In the U.S. bishops' commentary on the prayer, a Christian already comes out to meet Christ on the way to church. Christ meets the people in their gathering, in the word they hear and in the Communion they receive. He is also encountered in someone who needs their service. That last encounter may ask more of a person than coming to church has done. There is a reason the prayer asks for a willing response. One can see a neighbor in need and be thoroughly unwilling to become involved. When the priest prays for the will itself he includes that person in the petition.\footnote{USCCB, \textit{Commentary on the Proper Prayers of Advent from the Roman Missal}, First Sunday, Collect.}

\subsection{Israel's prayer and the answer the congregation sings}

The first name for God in Sunday's Isaiah passage is Father. Just before the reading begins there is an appeal beyond Abraham and Israel. The ancestors who gave this people its name can offer no present rescue. Earlier in the lament the speakers remembered how the people had come up from the sea and how God's Spirit had given them rest (Is 63:7--14). After the reading, they will describe the sanctuary gone up in flames (native 64:10--11). Calling God Father in these surroundings makes a claim about his relation to a people whose possessions, achievements and even place of worship can no longer speak in their favor.

In the U.S. reading the people ask God to descend, and the mountains quake before him. In Challoner the mountains melt at his presence (Is 64:1, 3). To ask God to act again as he once did is to admit that the lost people and their ruined city lie beyond anything they can put right by themselves. Then the object of their scrutiny changes. Looking at themselves they find a leaf loosened from its tree. The wind that has carried it away is their own iniquity. There is something different about that image from an account in which an enemy simply comes and carries them off. The power to which they have yielded has its origin in their own conduct. As the prayer continues, the failure even to call on God becomes one of the things they have to confess. The speakers finally resort to a different image from the natural world. Clay has no claim to have shaped itself. Each piece of work has a claim on the attention of the one who made it. Israel can still say whose people they are.

Different Bibles put the chapter break in different places. The passage assigned in the United States moves from a wish to be found doing right into an admission of guilt. A reader opening the Vulgate-based Challoner Bible will instead find a statement about God meeting someone who already does justice. That Bible also lets the sinful speakers go on to profess confidence in their future salvation before they speak of their uncleanness. Those features belong to the wording Jerome explains. The distinction affects whether one hears a request for a changed life or recalls someone in whom such a life has already been found. The two forms are identified beside the study passage.

In Psalm 80 the congregation sings about a vine that once took root. The rest of the psalm tells the story behind the selected verses. God brought the vine from Egypt and cleared the ground in which to plant it. In a verse outside Sunday's strophes, the psalmist asks God why he tore down its enclosure (Challoner Ps 79:13). Passersby can pluck its fruit and wild beasts can ravage it. The people ask the very one they hold responsible for the breach to return and tend his vine. The request to visit turns on the planter's enduring interest in his own work. A traveler could look at the damage and go his way. To the one who planted the vine, the ruined growth is the loss of something of his own.

The verses sung after Isaiah do not tell the Exodus story in full. They call to the shepherd of Israel, then turn to the vineyard, then speak about a man strengthened by God's right hand. At the end of the strophe the speakers promise what \emph{they} will do if given life. The changes of person matter. The people pray about someone before promising to call on God themselves. The response returns to the petition for God's face to shine. A member of the congregation who has lost the thread of the images can still join in asking for the favor of the one who sees them. Augustine recalls the humiliation of the crucified Christ when he considers how God's face could be present and yet scarcely discerned. To those who mocked the suffering man his weakness was visible enough. The psalmist's petition acquires a Christian use in the mouths of people who want to recognize him.\footnote{Psalm 80 (79), especially the vine in vv.~9--17. Augustine, \textit{Exposition on Psalm} 79, §§1--5, 9--11 (English heading Psalm 80).}

\subsection{A thanksgiving addressed to divided people}

Paul opens by wishing grace and peace to the Corinthians. The people addressed have already been called holy in the greeting immediately before the reading. In the lines immediately after it Paul asks them to stop quarreling. Between address and rebuke he describes a striking abundance of gifts. They can speak about the faith, they have received knowledge, and he can tell them that while they wait for Christ their church lacks no gift. A reader who shared their enthusiasm might expect a letter congratulating them on the health of their community. Paul gives thanks in such a way that he can go on to say exactly what must be corrected.

Paul thanks God for the Corinthians' speech and knowledge. God gave them grace in Christ. The testimony about Christ took root among them. It is Christ whom they expect to see revealed. A Corinthian proud of his ability to speak will receive a different impression of his achievement if he listens to all the people Paul credits for it. The names that mark parties later in the chapter, including Paul's own, are conspicuously ill suited to displace the one whose name recurs in this thanksgiving. Chrysostom pictures someone rousing drowsy hearers by bringing Christ's name to their attention again and again. By the time Paul asks whether he himself was crucified for them, they have already been reminded to whom they owe everything for which he praised them.\footnote{1 Corinthians 1:2--13, 18--31. John Chrysostom, \textit{Homilies on First Corinthians}, Homily 1, §§3--4, and Homily 2, §§1--8.}

The end of the passage leaves a further question. The Corinthians have received much and Paul expects God to strengthen them right through to the day of Christ. If the promise of future help meant their present conduct was satisfactory, his rebuke a few lines later would be hard to understand. If their divisions meant God had repudiated them, the rebuke would be hard to bear. Paul can appeal for change because the fellowship into which they were called still makes a claim on them. To break into groups devoted to rival teachers is to behave like people who have forgotten whom they have in common.

\subsection{A petition before the Gospel and a house left in trust}

At the Alleluia the congregation asks to see God's mercy. Psalm 85 has already recalled favors to Jacob. Why should a people whose wrongs have been forgiven still need to ask God for salvation? The psalm can both recall what God has done and look about for signs of the life that pardon makes possible. The singer expects an answer and presently falls silent to listen for it. Augustine hears the petition for salvation as a desire to see Christ himself. The congregation that makes the petition before the Gospel will stand and listen to what Christ says. Later it will receive him at Communion with another verse from this same psalm.\footnote{Augustine, \textit{Exposition on Psalm} 84, §§2--7 (English heading Psalm 85).}

Mark tells how Peter, James, John and Andrew asked Jesus privately about the fall of the Temple. His discourse includes terrible events, misleading announcements, persecution and the gathering of the elect. As he speaks of the fig tree in leaf, Jesus expects the disciples to recognize a sign just as they recognize the approach of summer. But no one hearing this discourse receives a date on which to cease watching. In the closing parable a master's time of arrival could fall in the evening or at midnight, at cockcrow or in the morning. Those possibilities cover a person's period of responsibility through the night. Nothing in the list makes any one of them the appointment.

A house with someone still at its door bears signs of an owner who means to come back. The operation of the house in his absence belongs to servants who exercise the authority he has given them. The Gospel does not specify what each servant makes or where each must stand. It does specify one service, that of the porter. The master took the trouble to tell him to keep watch. At the end Jesus tells everyone to do so. The difference between the porter's particular office and this command for all becomes clearer when we see how readily someone else in the story can fall asleep. Peter, James and John, three of the four who asked the question on the Mount of Olives, will go with Jesus into Gethsemane. Peter will make great promises about what he can endure, then be unable to stay awake with him for an hour. The disciples' eyes are heavy. A few pages separate the church's hearing of Jesus' warning from the account of his finding his friends asleep (Mark 14:29--42).

\subsection{Bringing gifts, receiving Communion, going home}

At the altar lie things that have passed through human hands. The USCCB commentator distinguishes the grain and grapes that grow from the bread and wine people have made from them. The distinction permits the worshippers to say quite truthfully that they have brought their gifts, even as they acknowledge that they first received them. There is a further human contribution at Mass when a person gives from what his work has earned so that someone else can eat. The gifts placed before God implicate their bearers. Someone who wished to make an offering while reserving for himself all claim over his own time and goods would have difficulty explaining where he obtained the things he can give. The Prayer over the Offerings also concerns redemption taking effect among those who offer. Its petition reaches into the life of people who can already receive benefits from what Christ has done.\footnote{USCCB, \textit{Commentary on the Proper Prayers of Advent}, First Sunday, Prayer over the Offerings.}

In Preface I of Advent the priest addresses the Father. The human condition the Son assumed at his first coming and the glory of his future appearing belong to the same Son. The first coming has already altered the prospects of those who await the second. They have been given a promise on which to venture the deeds for which they prayed at the Collect. Cyril of Jerusalem, teaching the Creed, attends to an astonishing reversal of roles. Once the Son stood silent under judgment. When he comes again it will be his place to judge. Remembering what he once endured changes what Christians know of the person for whom they wait. The disciples who failed him in the garden, and the churches to which their preaching went, have heard what he undertook for them before they have to see him on his seat of judgment.\footnote{Cyril of Jerusalem, \textit{Catechetical Lecture} XV, §§1--3, on the article of the Creed concerning the return of Christ in glory. U.S. Roman Missal, Preface I of Advent.}

At Communion one verse speaks of the Lord giving goodness and the earth yielding fruit. A hungry person depends on both. To come to the table is already to experience that dependence quite bodily. The two verbs in the verse allow for more than one question. What has God given to those who eat? What might he bring to maturity in their lives? In Augustine's exposition of this verse an acquired taste for justice shows itself in a changed use of money. Bellarmine sees Christ himself in the earth's fruit. Even a reader familiar with both themes elsewhere in the psalm needs to distinguish what each author says of \emph{this} verse, for it is the one the Missal appoints. Their accounts are developed below.

The Prayer after Communion concerns the life in which communicants will find themselves again when Mass ends. The USCCB commentary dwells on what the worshippers have just done with the products of the earth. Real bread, something made and consumed, has been brought to the altar. In sacramental Communion the people receive Christ's Body. Something that daily labor provides can enter a person's life as the means of a gift far beyond the power of that labor to produce. Learning to love goods in the right order involves remembering both who gave them and what they can be used for. The lesson has applications which the giver of a meal, the person dependent on a wage and the recipient of someone else's provision will have different opportunities to discover.\footnote{USCCB, \textit{Commentary on the Proper Prayers of Advent}, First Sunday, Prayer after Communion.}

If the priest uses the solemn blessing for Advent, the assembly responds to three invocations before the final blessing. The first looks back to the coming that has taken place. The last looks ahead to the meeting still to come. Between them the priest asks for faith, hope and charity in the race the people are running now. Those who say Amen can go from the church with something more definite to ask of the days ahead. They will have occasion to believe someone whose face they cannot yet see, to keep their place in the house he has entrusted to them, and to exercise their charity toward someone whose need they can already see.
